\section{Where AR Falls}
\label{sec:ar-position}

\subsection{Estimate: Within One Seat of Ensemble Median for 5/6 States}

Using the G.1 results, we summarise the \ar{} partisan position:

\begin{center}
\begin{tabular}{lcccc}
\toprule
State & $k$ & AR seats (D) & Ensemble median (D) & AR percentile \\
\midrule
NC & 14 & 7  & 7    & 54th \\
WI &  8 & 4  & 4    & 61st \\
GA & 14 & 5  & 6    & 38th \\
PA & 17 & 8  & 8    & 52nd \\
TX & 38 & 15 & 15\est & 52nd\est \\
CA & 52 & 38 & 37\est & 55th\est \\
\midrule
\multicolumn{4}{l}{\textbf{States within 1 seat of median}} & 5/6 \\
\bottomrule
\multicolumn{5}{l}{\small \est~Literature-bound estimate.}
\end{tabular}
\end{center}

\paragraph{Interpretation.}
For NC, WI, and PA, the \ar{} plan exactly matches the ensemble median.
For TX and CA, the \ar{} plan is within one seat of the estimated median.
Only Georgia diverges meaningfully: the \ar{} plan produces one fewer
Democratic seat than the ensemble median of 6.

\subsection{The Georgia Deviation}

The Georgia result warrants detailed examination.  With $k = 14 = 2 \times 7$,
the \ar{} plan produces a binary top-level split of the state, then a 7-way
partition of each half.  The binary split divides the state roughly along a
line from northeast to southwest, separating the Atlanta metropolitan area
from the western and southern parts of the state.

The Atlanta half contains the Democratic concentration (Fulton, DeKalb,
Gwinnett, Clayton counties plus the Black Belt corridor).  Under a 7-way
partition of the Atlanta half, the algorithm produces 5 Democratic-leaning
districts and 2 Republican-leaning suburban districts.  The remaining half
produces 7 Republican-leaning districts.

Ensemble plans with 6 Democratic seats typically achieve this by using a
less compact boundary that captures more of the Atlanta suburbs in the
Democratic half.  But such a boundary has a higher edge cut and is therefore
disfavored by the METIS objective.  In other words, the \ar{} plan trades
one Democratic seat for a lower edge cut.  This is a feature of the algorithm,
not a partisan choice.

\subsection{Robustness to Electoral Baseline}

The partisan percentiles reported above use the 2020 presidential vote as the
baseline.  We also computed percentiles under the 2018 midterm baseline
(more Democratic-leaning) and the 2016 presidential baseline (slightly less
Democratic-leaning).  The \ar{} percentile rankings are stable across
baselines to within 5 percentage points, confirming that the results are not
sensitive to baseline choice.
